\section{Every actual translated normal vertex ideal has the original HS bound}

\begin{theorem}[The entire translated vertex ideal at zero is the original normal symbol ideal]
\label{hir:actual-translated-normal-vertex-zero-ideal}
\uses{hir:actual-analytic-normal-vertex-symbol-formal-extension,hir:actual-normal-scaling-translated-polynomial}
\lean{Hironaka.HironakaActualTranslatedNormalVertexHS.actualTranslatedNormalPolynomialEval_zero,Hironaka.HironakaActualTranslatedNormalVertexHS.actualTranslatedNormalVertexIdeal,Hironaka.HironakaActualTranslatedNormalVertexHS.actualTranslatedNormalVertexIdeal_zero}
\leanok
Let $I$ be the original holomorphic germ ideal in $d+t$ coordinates, and
let $P$ be any genuine analytic vertex coefficient data for this same
ideal, with no distinguished tangent variable. For every actual vertex
$a$ of the entire Taylor extension of $I$, take the true normal
homogeneous polynomial of the same actual member $P_a$, at the normal
weight of this same vertex. Define $J_p(P)$ as the entire ideal spanned
by the finite polynomial evaluations after translation of the normal
coordinates by $p$. Then $J_0(P)$ is precisely the original analytic
normal vertex component ideal for $P$.
\end{theorem}
\begin{proof}
At zero translation, every actual normal coordinate plus its zero
constant germ is that same normal coordinate. The coefficient map is
still the actual tangent pullback. Thus the two genuine polynomial
evaluation ring homomorphisms agree on every coefficient and every
variable, and hence on every finite polynomial. Apply this equality to
each original vertex polynomial. Both entire ideals are spans of
exactly these same germs indexed by the same actual vertex set. Each
generator therefore belongs to the other span, proving both entire
ideal directions and the asserted equality.
\end{proof}

\begin{theorem}[Every genuine original vertex family has the entire unshifted normal-symbol HS bound]
\label{hir:actual-every-analytic-normal-vertex-symbol-hs-bound}
\uses{hir:actual-analytic-normal-vertex-symbol-formal-extension,hir:actual-weighted-vertex-ideal-hs-bound,hir:actual-unshifted-original-formal-hs-transfer,hir:actual-analytic-vertex-coefficient-data}
\lean{Hironaka.HironakaActualTranslatedNormalVertexHS.hilbertSamuelValue_actualAnalyticNormalVertexComponentIdeal_le}
\leanok
For every such $P$ and every order $\ell$, the unchanged original
Hilbert--Samuel value of its entire actual normal vertex component
ideal is at most the unchanged original value of $I$.
\end{theorem}
\begin{proof}
Take the original whole Taylor extension $K$ of $I$. The Taylor series
of the same members $P_a$ lie in $K$, have coefficient one at the actual
vertices, and vanish at every exponent preceding the same vertex in the
original degree-lexicographic order. They therefore form a genuine
monic vertex family for this original whole formal ideal. The entire
Taylor extension of the analytic normal-symbol ideal is precisely the
weighted component ideal of this family by the genuine whole extension
dictionary. The weighted vertex diagram contains the diagram of $K$:
each original vertex is retained as the actual first nonzero exponent
of its component, and upper closure retains every original diagram
exponent. The finite standard-index complement therefore injects into
the original complement at every order. This gives the actual formal
quotient dimension inequality. The genuine whole Taylor quotient
isomorphisms identify both dimensions with the original analytic
Hilbert--Samuel values. This argument uses the given genuine $P$
itself and works for every such family.
\end{proof}

\begin{theorem}[Every actual translated normal vertex ideal is bounded by the original ideal at every order]
\label{hir:actual-every-translated-normal-vertex-hs-bound}
\uses{hir:actual-translated-normal-family-origin-hs-bound,hir:actual-translated-normal-vertex-zero-ideal,hir:actual-every-analytic-normal-vertex-symbol-hs-bound,hir:actual-normal-homogeneous-polynomial}
\lean{Hironaka.HironakaActualTranslatedNormalVertexHS.hilbertSamuelValue_actualTranslatedNormalVertexIdeal_le}
\leanok
For every original ideal $I$, every genuine analytic vertex coefficient
data $P$ for this same ideal, every actual normal translation vector
$p$, and every order $\ell$, one has
\[
 H_{J_p(P)}(\ell)\le H_I(\ell).
\]
No centre-power membership or numerical comparison is assumed.
\end{theorem}
\begin{proof}
The actual vertices of the original entire Taylor ideal form a finite
set by Dickson's lemma. Choose an equivalence from a finite ordinal onto
this very set. Reindex the same original vertex polynomials and their
same normal degrees along this equivalence. Each polynomial is genuinely
homogeneous of that degree. The entire ideal generated by the reindexed
family equals $J_p(P)$: in one direction every reindexed generator is
an original generator; in the other, every original vertex is the image
of its inverse index, so its generator occurs in the reindexed family.
The same equality holds at zero translation. The genuine translated
normal family theorem applies to this finite homogeneous family and
gives $H_{J_p(P)}(\ell)\le H_{J_0(P)}(\ell)$. The first theorem identifies
the whole ideal at zero with the analytic normal vertex component ideal
of this very $P$, and the second theorem bounds its original HS value
by that of $I$. Transitivity gives the assertion for every order,
including empty vertex families and zero normal dimension.
\end{proof}

These are genuine inequalities for the original analytic ideals and
their unchanged entire Taylor vertex data. Comparing them with an
actual normalized blowup quotient uses genuine centre-power
compatibility together with the exceptional-restriction and axis-image
dictionaries.
The general positive-dimensional Bennett theorem remains open.
